\documentclass[11pt]{article}
\usepackage[T1]{fontenc}
\usepackage{lmodern}
\usepackage[margin=1in]{geometry}
\usepackage{amsmath,amssymb,amsthm}
\usepackage{xurl}
\usepackage[hidelinks]{hyperref}
\hypersetup{pdftitle={A degree-three counterexample for the normalized upper Laplacian},pdfauthor={Kaiwen Jin}}
\newtheorem{theorem}{Theorem}
\newtheorem{lemma}[theorem]{Lemma}
\newcommand{\Spec}{\operatorname{Spec}}
\newcommand{\degf}{\operatorname{deg}}
\newcommand{\norm}[1]{\left\lVert#1\right\rVert}
\newcommand{\ip}[2]{\langle#1,#2\rangle}
\begin{document}

\title{A degree-three counterexample for the normalized upper Laplacian}
\author{Kaiwen Jin\\\small\href{mailto:yui.ui.siki@gmail.com}{\texttt{yui.ui.siki@gmail.com}}}
\date{14 October 2026}
\maketitle

\begin{abstract}
Does a low codimension-one degree force zero into the spectrum of the
normalized upper Laplacian of an infinite simplicial complex? We construct a
countable pure three-dimensional complex with maximum triangle degree three
whose upper Laplacian has spectrum bounded away from zero on the full
degree-weighted triangle cochain space. This disproves both assertions in
AIM Problem 643, including the stronger claim that degree at most the
dimension forces a nonzero square-summable kernel. Every edge of the example
has infinite triangle incidence. The proof uses disjoint square-summable
fillings to control selected triangle values, then the gradient estimate on
a binary tree to control all remaining coordinates.
\end{abstract}

\section{The endpoint questions}\label{sec:introduction}

The upper simplicial Laplacian measures the failure of a cochain to be a
cocycle. Its normalized form also arises in the expected signed particle
counts of simplicial branching random walks. Rosenthal computes the spectral
measure on regular arboreal complexes and asks whether the corresponding
endpoint eigenvalue persists for general low-degree complexes
\cite[Theorem 1.1 and Conjecture 4.3]{Rosenthal}.

We use the precise formulation in AIM Problem 643~\cite{AIM643}.
For a nonempty pure $d$-dimensional simplicial complex $X$, with $d\geq2$,
let $\degf(F)$ be the number of $d$-simplices containing a $(d-1)$-simplex
$F$, and set $D=\sup_F\degf(F)$. The full cochain spaces are
\[
 H=\ell^2(X^{d-1},\degf),\qquad K=\ell^2(X^d),\qquad
 \norm{f}_H^2=\sum_F\degf(F)|f(F)|^2,
\]
using one orientation per simplex and counting measure in $K$. The usual
coboundary is
\[
 (\delta f)([x_0,\ldots,x_d])
   =\sum_{i=0}^d(-1)^if([x_0,\ldots,\widehat x_i,\ldots,x_d]),
 \qquad L=\delta^*\delta.
\]
The two assertions are that $D\leq d+1$ forces $0\in\Spec(L)$, and that
$D\leq d$ forces $\ker L\ne\{0\}$. The complex may be infinite, and no
finite-incidence condition is imposed on lower-dimensional faces.

\begin{theorem}\label{thm:main}
There exists a countable nonempty pure three-dimensional simplicial complex
$X$ with one triangle of degree one and every other triangle of degree three.
Every edge belongs to infinitely many triangles. For its full spaces
$H=\ell^2(X^2,\degf)$ and $K=\ell^2(X^3)$, the usual coboundary satisfies
\[
 \norm{f}_H^2\leq C\norm{\delta f}_K^2\quad(f\in H),\qquad
 C=12+\frac{3(1+2\sqrt3)^2}{(\sqrt2-1)^2}.
\]
Consequently
\[
 \Spec(\delta^*\delta)\subset[C^{-1},4],\qquad
 \ker(\delta^*\delta)=\{0\}.
\]
\end{theorem}

Here $D=3\leq d+1=4$, while zero is absent from the spectrum. Also
$D=3\leq d=3$, while the square-summable kernel is zero. Thus the same
example disproves each universal assertion separately.

The construction repeatedly attaches copies of an apex triangle and a
binary tree. A unit flux splits equally at each tree vertex, giving a
square-summable filling of the apex triangle. The fresh vertices make these
fillings orthogonal across copies. Their pairings with $\delta f$ control
the triangle values used as attachment targets. The body coboundary equations
then express the remaining triangle values through tree gradients, which
have a uniform lower bound. This last step gives control of the full cochain
space rather than only the target coordinates.

\section{Construction and face degrees}\label{sec:construction}

Let $W=\{0,1\}^{<\omega}$ be the finite binary words, including the empty
word $o$. The rooted binary tree $G$ has oriented edges $w\to ws$, where
$s\in\{0,1\}$. Its root has degree two and every other vertex has degree
three; every vertex has two children. Write $|w|$ for word length.

We first describe one part of the construction. Given an ordered triangle
$A=(a,b,c)$ and a fresh copy of $G$, with vertices $v_w$, insert
\begin{equation}\label{eq:tetrahedra}
 [a,b,v_w,v_{ws}],\quad [b,c,v_w,v_{ws}],\quad
 [c,a,v_w,v_{ws}]
 \qquad(w\in W,\ s\in\{0,1\}),
\end{equation}
and the single tetrahedron $[a,b,c,v_o]$. The first three families will be
called body tetrahedra and the last one the bridge. Each triangle
$[p,v_w,v_{ws}]$, with $p\in\{a,b,c\}$, is then used as the ordered apex
triangle of a new copy of this construction. This is repeated at every
generation.

Here is a precise vertex labeling for the recursion. Set
\[
 \mathcal B=\{0,1,2\}\times W\times\{0,1\},
 \qquad \mathcal I=\mathcal B^{<\omega}.
\]
For every $J\in\mathcal I$ and $w\in W$, take a distinct vertex $v_{J,w}$.
Take three further vertices $r_0,r_1,r_2$, with different tags. Define ordered
apex triples recursively by
\begin{equation}\label{eq:recursion}
 A_{\varnothing}=(r_0,r_1,r_2),\qquad
 A_{J*(p,w,s)}=(a_{J,p},v_{J,w},v_{J,ws}),
 \quad A_J=(a_{J,0},a_{J,1},a_{J,2}).
\end{equation}
The apex vertices of $J$ come from smaller generations, so they are distinct
from its fresh vertices $v_{J,w}$; the three apex vertices are distinct by
induction. Let $X$ be the downward closure of the tetrahedra in
\eqref{eq:tetrahedra}, using $A_J$ and $v_{J,w}$ for every $J$, together with
their bridges. No higher-dimensional simplices are added. Thus $X$ is a
countable, nonempty, pure three-dimensional abstract simplicial complex.

The triangles first appearing in copy $J$ have two forms:
\begin{align}
 E(J,p,w,s)&=[a_{J,p},v_{J,w},v_{J,ws}],\label{eq:E}\\
 V(J,pq,w)&=[a_{J,p},a_{J,q},v_{J,w}],
       \qquad pq\in\{01,12,20\}.\label{eq:V}
\end{align}
The bridge has the additional triangle $F_J=[A_J]$. For $J=\varnothing$ this
is $F_0=[r_0,r_1,r_2]$; otherwise it is the $E$ triangle specified by the last
letter of $J$. We use \eqref{eq:E}--\eqref{eq:V} as global triangle
orientations. An inherited target $F_J$ keeps its existing orientation, as
prescribed by \eqref{eq:recursion}. The displayed tetrahedra fix their own
orientations. Other orderings have the usual alternating sign.

\begin{lemma}\label{lem:degrees}
Every triangle of $X$ has a unique first appearance as $F_0$, $E$, or $V$.
The tetrahedron sets of different copies are disjoint. Each $E$ triangle is
the target $F_J$ of exactly one child copy. The triangle $F_0$ has degree one
and all $E$ and $V$ triangles have degree three. Every edge of $X$ belongs to
infinitely many triangles.
\end{lemma}
\begin{proof}
A new $E$ or $V$ triangle contains a vertex with the current copy label $J$.
No smaller-generation copy contains that vertex. New faces in a different
copy of the same generation have its different label, while their apex
vertices come from smaller generations. New faces at later generations
contain later labels. Within one copy, an $E$ triangle has two new vertices
and a $V$ triangle has one; the apex vertex or edge and tree coordinates
distinguish them. The only old triangle in a copy is its bridge target.
The same argument distinguishes tetrahedra: bodies have two new vertices,
bridges have one, and different tree edges and apex edges distinguish bodies.
Formula \eqref{eq:recursion} therefore gives a bijection between $E$
triangles and noninitial copy targets. No $V$ triangle is targeted.

An $E$ triangle has two body cofaces in its birth copy and one bridge coface
in its child. A $V$ triangle at the tree root has two body cofaces and the
bridge; away from the root it has three body cofaces. These are all cofaces,
since a later tetrahedron can contain an old triangle only as its bridge
target. The initial triangle has only the initial bridge coface.

Every apex edge has the distinct $V$ cofaces indexed by all $w\in W$.
A tree edge, or an edge joining an apex vertex to a tree vertex, lies in an
$E$ triangle and hence becomes an apex edge in the corresponding child.
These exhaust the edge types in $X$, proving the last assertion.
\end{proof}

\section{Square-summable fillings}\label{sec:fillings}

Put $m(F)=\degf(F)$. Let $H$ be the full space of complex triangle cochains
with degree-weighted norm, and let $K$ be the counting $\ell^2$ space of
tetrahedron cochains. Lemma~\ref{lem:degrees} gives
\begin{equation}\label{eq:Hpartition}
 \norm{f}_H^2=|f(F_0)|^2+3\sum_E|f(E)|^2+3\sum_V|f(V)|^2.
\end{equation}
All sums use one chosen orientation per simplex. Inner products are linear
in the first variable. On finitely supported cochains, Cauchy--Schwarz on
the four facets of a tetrahedron gives
\begin{equation}\label{eq:bounded}
 \norm{\delta f}_K^2\leq4\sum_Fm(F)|f(F)|^2=4\norm{f}_H^2.
\end{equation}
Consequently $\delta$ extends to all $H$ with norm at most two. Continuous
coordinate evaluations show that its extension is still the usual
alternating facet sum. Its bounded adjoint is
\begin{equation}\label{eq:adjoint}
 (\delta^*h)(F)=\frac1{m(F)}
     \sum_{T\supset F}\epsilon(T,F)h(T),
\end{equation}
where $\epsilon(T,F)$ is the boundary incidence sign. Each sum has at most
three terms. To verify membership in $H$ without any lower-face assumption,
pair the right side with finitely supported $f$. The finite rearrangement
equals $\ip{\delta f}{h}_K$ and is bounded by
$2\norm{f}_H\norm{h}_K$. Taking $f$ to be finite restrictions of that right
side shows its $H$ norm is at most $2\norm{h}_K$; density then proves
\eqref{eq:adjoint}. In particular $L=\delta^*\delta$ is bounded,
self-adjoint and positive, with $\norm{L}\leq4$.

We next fill the target $F_J$ in its own copy. Fix $J$ temporarily, suppress
its label, and write $A=[a,b,c]$ and
$z=[a,b]+[b,c]+[c,a]$. For $N\geq0$, let $c_N$ be the finite tree
one-chain having coefficient $2^{-n}$ on every edge whose child has depth
$n\leq N$. With counting norms on graph edges,
\begin{equation}\label{eq:treeflow}
 \partial c_N=-[v_o]+2^{-N}\sum_{|w|=N}[v_w],
 \qquad \norm{c_N}^2=\sum_{n=1}^N2^{-n}=1-2^{-N}.
\end{equation}
This includes $c_0=0$. At an interior vertex the incoming coefficient equals
the sum of the two outgoing coefficients.

For finite chains, use $*$ to concatenate ordered vertices, extending
bilinearly. Directly,
\[
 \partial[p,q,u,v]=[q,u,v]-[p,u,v]+[p,q,v]-[p,q,u].
\]
Summing over the three apex cycle edges cancels the first two terms. Also
$\partial[a,b,c,v_o]=z*[v_o]-A$. Therefore the finite top chain
$h_{J,N}=-z*c_N-A*[v_o]$, identified with a cochain in $K$, satisfies
\begin{equation}\label{eq:finitefilling}
 \partial h_{J,N}=F_J-2^{-N}z*\sum_{|w|=N}[v_w],
 \qquad \norm{h_{J,N}}_K^2=4-3\,2^{-N}.
\end{equation}
Only finite formal boundaries have been used. The tail squared norm after
depth $N$ is $3\,2^{-N}$, so $h_{J,N}$ converges strongly in $K$ to a
cochain $h_J$ with $\norm{h_J}_K^2=4$.

The $3\cdot2^N$ leaf triangles in \eqref{eq:finitefilling} are distinct $V$
triangles of actual final degree three, also when $N=0$. If $\mathbf1_F$
denotes the cochain supported at $F$ with value one, \eqref{eq:adjoint} gives
\[
 \delta^*h_{J,N}=\frac{\mathbf1_{F_J}}{m(F_J)}+r_{J,N},
 \qquad
 \norm{r_{J,N}}_H^2
   =(3\cdot2^N)\,3\left(\frac{2^{-N}}3\right)^2=2^{-N}.
\]
The boundedness of $\delta^*$ now proves
\begin{equation}\label{eq:fillinglimit}
 \delta^*h_J=\frac{\mathbf1_{F_J}}{m(F_J)},\qquad
 \norm{h_J}_K^2=4.
\end{equation}
Thus no formal boundary at an infinitely incident edge is needed.

\section{A lower bound on the full cochain space}\label{sec:estimate}

The cochains $h_J$ have disjoint tetrahedron supports by
Lemma~\ref{lem:degrees}. Hence $h_J/2$ form an orthonormal family in $K$.
For every $f\in H$, \eqref{eq:fillinglimit} gives
\[
 \ip{\delta f}{h_J}_K
   =\ip{f}{\mathbf1_{F_J}/m(F_J)}_H=f(F_J).
\]
Subtracting the orthogonal projection onto any finite subfamily from
$\delta f$ and expanding its squared norm gives the finite Bessel inequality.
Increasing those finite sets to all $J$, and recalling that the targets are
exactly $F_0$ and all $E$ triangles, yields
\begin{align}
 |f(F_0)|^2+\sum_E|f(E)|^2&\leq4\norm{\delta f}_K^2,
       \label{eq:targets}\\
 |f(F_0)|^2+3\sum_E|f(E)|^2&\leq12\norm{\delta f}_K^2.
       \label{eq:weightedtargets}
\end{align}

It remains to control the $V$ coordinates. Within each copy, pull tree
coordinates back along $w\mapsto v_{J,w}$ and write
\[
 e_p(w,ws)=f([a_{J,p},v_{J,w},v_{J,ws}]),\qquad
 g_{pq}(w)=f([a_{J,p},a_{J,q},v_{J,w}]),
 \quad pq\in\{01,12,20\}.
\]
Let $e$ and $g$ denote their counting $\ell^2$ direct sums over all copies
and indicated apex vertices or edges. The unique first appearances of the
faces ensure that these are precisely the $E$ and $V$ coordinates, without
duplication. In particular \eqref{eq:targets} gives
$\norm{e}\leq2\norm{\delta f}_K$.

Let $y$ be $\delta f$ restricted to body tetrahedra, so
$\norm{y}\leq\norm{\delta f}_K$. Their alternating boundaries give
\begin{equation}\label{eq:bodyidentity}
 y_{pq}(w,ws)=g_{pq}(ws)-g_{pq}(w)+e_q(w,ws)-e_p(w,ws).
\end{equation}
For each tree edge, the last three apex differences are the matrix
\[
 B=\begin{pmatrix}-1&1&0\\0&-1&1\\1&0&-1\end{pmatrix}
\]
applied to $(e_0,e_1,e_2)$. Multiplication gives
$B^*B=3I-\mathbf{J}_3$, where $\mathbf{J}_3$ is the all-ones matrix. Thus
$\norm{Bv}^2=3\norm{v}^2-|v_0+v_1+v_2|^2\leq3\norm{v}^2$.
Taking the direct sum of these matrices in \eqref{eq:bodyidentity}, we obtain
\begin{equation}\label{eq:gradientupper}
 \norm{\nabla g}\leq\norm{y}+\norm{Be}
    \leq(1+2\sqrt3)\norm{\delta f}_K,
 \qquad (\nabla g)_{pq}(w,ws)=g_{pq}(ws)-g_{pq}(w).
\end{equation}

For a scalar $\ell^2(W)$ function $q$, index tree edges by their child and
let $Sq$ and $Tq$ be the child and parent values respectively. Then
\[
 \norm{Sq}^2=\norm{q}^2-|q(o)|^2,\qquad
 \norm{Tq}^2=2\norm{q}^2.
\]
The reverse triangle inequality implies
$\norm{Sq-Tq}\geq(\sqrt2-1)\norm{q}$. The same identities hold on the
direct sum over all copies and apex edges, with the root terms summed. Hence
\begin{equation}\label{eq:Vbound}
 3\sum_V|f(V)|^2=3\norm{g}^2
 \leq\frac{3(1+2\sqrt3)^2}{(\sqrt2-1)^2}\norm{\delta f}_K^2.
\end{equation}
Combining \eqref{eq:Hpartition}, \eqref{eq:weightedtargets} and
\eqref{eq:Vbound} proves
\begin{equation}\label{eq:coercivity}
 \norm{f}_H^2\leq C\norm{\delta f}_K^2\quad\text{for every }f\in H,
 \qquad C=12+\frac{3(1+2\sqrt3)^2}{(\sqrt2-1)^2}.
\end{equation}
All direct-sum maps above are bounded and all estimates use the actual final
face degrees. In particular the estimate applies to the whole $H$.

\begin{proof}[Completion of the proof of Theorem~\ref{thm:main}]
Put $\alpha=C^{-1}$. Equations \eqref{eq:bounded} and \eqref{eq:coercivity}
give
\begin{equation}\label{eq:formbounds}
 \alpha\norm{f}_H^2\leq\ip{Lf}{f}_H\leq4\norm{f}_H^2.
\end{equation}
For completeness we verify the resolvent conclusion directly. A bounded
Hilbert-space operator $T$ satisfying positive lower bounds for both $T$ and
$T^*$ is invertible: a convergent sequence $Tx_n$ makes $x_n$ Cauchy, so the
range is closed, and $(\operatorname{ran}T)^\perp=\ker T^*=\{0\}$.
A closed subspace with zero orthogonal complement is the whole space. Indeed,
for a proper closed subspace $M$ and $v\notin M$, let
$\rho=\inf_{m\in M}\norm{v-m}>0$ and choose a minimizing sequence $m_n$,
with $a_n=\norm{v-m_n}\to\rho$. The parallelogram identity gives
\[
 \norm{m_n-m_\ell}^2
 \leq2(a_n^2-\rho^2)+2(a_\ell^2-\rho^2).
\]
It therefore converges to $m\in M$ with $\norm{v-m}=\rho$. Varying $m$ by
real and imaginary multiples of a vector in $M$ shows that $v-m$ is a
nonzero vector in $M^\perp$. The lower bound for $T$ also bounds its inverse.

For nonreal $\lambda$, the imaginary part of
$\ip{(L-\lambda I)f}{f}_H$ gives a lower bound
$|\operatorname{Im}\lambda|\norm{f}_H$ for
$\norm{(L-\lambda I)f}_H$, and likewise for its adjoint
$L-\overline\lambda I$. For real $\lambda<\alpha$, \eqref{eq:formbounds}
gives the lower bound $\alpha-\lambda$ for $L-\lambda I$ and its adjoint;
for $\lambda>4$ it gives the bound $\lambda-4$ for $\lambda I-L$.
The criterion just proved shows that all such $\lambda$ are resolvent
points. Thus $\Spec(L)\subset[\alpha,4]$ and $\ker L=\{0\}$.
Lemma~\ref{lem:degrees} supplies all remaining geometric assertions.
\end{proof}

\begin{thebibliography}{2}
\bibitem{AIM643}
MColbrook/AIM problem repository,
\emph{643. Zero in the upper-Laplacian spectrum of low-degree simplicial
complexes}, statement at revision \texttt{a153f469},
\url{https://github.com/MColbrook/AIM/blob/a153f469ee0ea64e876b2ad97c79d252c19a3e22/problems/643-low-degree-simplicial-zero-spectrum.md}.
\bibitem{Rosenthal}
R. Rosenthal,
\emph{Simplicial branching random walks and their applications},
arXiv:1412.5406v1 (2014),
\url{https://arxiv.org/abs/1412.5406v1}.
\end{thebibliography}

\paragraph{AI use.}
The construction, proof development and manuscript preparation used
OpenAI Codex. Natural-language checking was also AI-assisted. Independent
mathematical review is requested.

\end{document}
